\chapter{Anti-Pattern Catalogue}
\label{app:ap-catalogue}

This appendix collects all CY-specific anti-patterns (AP-CY1 through
AP-CY33), cross-programme anti-patterns (AP150--AP157), and the
formula-mechanical pattern FM24 into a comprehensive reference.
Each entry records the failure mode, its severity, and the
counter-measure.  These patterns were identified through systematic
error archaeology across 180+ commits and a 180-agent verification
swarm.

%% =====================================================================
\section{Severity levels}
\label{sec:ap-sev}

\begin{center}
\begin{tabular}{@{}lp{4.5cm}p{5cm}@{}}
\toprule
\textbf{Level} & \textbf{Meaning} & \textbf{Required action} \\
\midrule
\textsc{Critical} & Theorem status wrong; logical
  error propagates to downstream results &
  Immediate fix; audit all instances in all three volumes \\
\textsc{High} & Numerical or structural error;
  formula incorrect; dimension wrong &
  Fix before next build \\
\textsc{Medium} & Convention clash, ambiguity,
  or scope confusion &
  Fix in current session \\
\textsc{Low} & Cosmetic, cross-reference staleness,
  or README-only issue &
  Fix in batch \\
\bottomrule
\end{tabular}
\end{center}

\noindent
\textbf{Statistics}:
\begin{center}
\begin{tabular}{@{}lr@{}}
\toprule
\textbf{Category} & \textbf{Count} \\
\midrule
CY-specific (AP-CY1--AP-CY52) & 52 \\
Cross-programme (AP150--AP157) & 8 \\
Formula-mechanical (FM24) & 1 \\
\midrule
\textbf{Total catalogued} & \textbf{61} \\
\midrule
Critical severity & 14 \\
High severity & 28 \\
Medium severity & 16 \\
Low severity & 1 \\
\bottomrule
\end{tabular}
\end{center}

\noindent
Two anti-patterns have required fixes in 10+ independent instances:
AP-CY6/AP-CY14 (unconstructed object in theorem environment,
11+ fixes) and AP113 (bare $\kappa$, 15+ fixes).  These two
alone account for approximately 40\% of all error-correction
commits in Vol~III.

%% =====================================================================
\section{CY-specific anti-patterns: AP-CY1 through AP-CY8}
\label{sec:ap-cy1-8-full}

{%
\small
\begin{longtable}{@{}p{1.6cm}p{2.4cm}lp{7.4cm}@{}}
\toprule
\textbf{ID} & \textbf{Name} & \textbf{Severity}
  & \textbf{Description and counter} \\
\midrule
\endhead
AP-CY1 & CY dim $\neq$ cpx dim & \textsc{High} &
  $\Fuk(X)$ and $D^b(\Coh(X))$ are CY$_n$ where $n$ is the
  \emph{complex} dimension, not the real dimension $2n$.
  \textbf{Counter}: always state ``CY$_d$ with
  $d = \dim_\C X$.'' \\
\midrule
AP-CY2 & CY trace target & \textsc{High} &
  The CY trace lives in $\HC^-_d(\cC)$ (negative cyclic homology),
  not just $\HH_d \to k$.  The negative cyclic refinement is
  essential for the $\Sd$-framing.
  \textbf{Counter}: always write $\HC^-_d$, never bare
  $\HH_d \to k$. \\
\midrule
AP-CY3 & $\Etwo \neq$ commutative & \textsc{High} &
  $\Etwo$ braiding is \emph{not} symmetric.
  $\Etwo \to \Einf$ loses all quantum group structure.
  \textbf{Counter}: write ``braided,'' never ``commutative''
  for $\Etwo$. \\
\midrule
AP-CY4 & Drinfeld $\neq$ derived ctr & \textsc{High} &
  $Z(\cC)$ (monoidal center via half-braidings) $\neq$
  $\Zder(A)$ (Hochschild cochains).
  Drinfeld center categorifies the derived center.
  \textbf{Counter}: always specify which center
  (see Appendix~\ref{app:notation-conventions},
  Section~\ref{sec:nc-centers}). \\
\midrule
AP-CY5 & Root-of-unity req. & \textsc{Medium} &
  Kazhdan--Lusztig equivalence requires $q$ a root of unity.
  At generic $q$, $\Rep_q(\frakg)$ is semisimple.
  \textbf{Counter}: state the $q$-specialization explicitly. \\
\midrule
AP-CY6 & $A_X$ at $d=3$ & \textsc{Critical} &
  $A_X$ for CY3 does \emph{not} exist---it IS the $d=3$
  programme.  Results depending on $A_X$ at $d=3$ must use
  \verb|\begin{conjecture}| and \verb|\ClaimStatusConditional|,
  naming CY-A$_3$.
  \textbf{Counter}: apply the decision tree (HZ3-1).
  Q1: proof chain through $A_X$ at $d=3$?
  $\to$ \verb|\begin{conjecture}|.
  Q2: through $A_X$ at $d=2$?
  $\to$ \verb|\begin{theorem}|, cite CY-A.
  Q3: pure categorical / VOA?
  $\to$ \verb|\begin{theorem}|, classical proof. \\
\midrule
AP-CY7 & CoHA $\neq$ $\Eone$-chiral & \textsc{High} &
  CoHA is associative (Schiffmann--Vasserot--Yang--Zhao
  Hall product), not a chiral algebra.
  ``$\Eone$-sector of $G(X)$'' assumes $G(X)$ exists (AP43).
  \textbf{Counter}: connection is via the functor $\Phi$,
  not by identification. \\
\midrule
AP-CY8 & Borcherds $\neq$ bar Euler & \textsc{High} &
  $\Phi_{10} =$ bar Euler product is an \emph{observation}
  (for K3$\times E$), not a theorem.
  Conditional on CY-A$_2$ and Vol~I Borcherds-lift
  identification.
  \textbf{Counter}: cite both CY-A and the Vol~I anchor
  explicitly. \\
\bottomrule
\end{longtable}
}%

%% =====================================================================
\section{Empirical anti-patterns: AP-CY9 through AP-CY13}
\label{sec:ap-cy9-13-full}

Identified through 50-commit error archaeology.

{%
\small
\begin{longtable}{@{}p{1.6cm}p{2.4cm}lp{7.4cm}@{}}
\toprule
\textbf{ID} & \textbf{Name} & \textbf{Severity}
  & \textbf{Description and counter} \\
\midrule
\endhead
AP-CY9 & Jacobi discriminant & \textsc{High} &
  For $\phi_{k,m}$ of index $m$, only discriminants $D$ with
  $D \equiv 0$ or $3 \pmod{4}$ (at $m=1$) can appear.
  $c(-1) = 2$ for $\phi_{0,1}$ in Eichler--Zagier convention,
  not $1$.
  \textbf{Counter}: verify discriminant constraint before
  filling coefficient tables. \\
\midrule
AP-CY10 & Flop $\neq$ Koszul dual & \textsc{High} &
  Birational flop $X \dashrightarrow X^+$ \emph{preserves}
  $\kappa_{\mathrm{ch}}$: $\kappa_{\mathrm{ch}}(A_X) =
  \kappa_{\mathrm{ch}}(A_{X^+})$.  Koszul dual $A^!$
  satisfies $\kappa_{\mathrm{ch}}(A) +
  \kappa_{\mathrm{ch}}(A^!) = \rho_K$.  Flop exchanges
  K\"{a}hler chambers; Koszul exchanges algebra/coalgebra.
  \textbf{Counter}: never write ``flop is the Koszul dual.'' \\
\midrule
AP-CY11 & Conditional transitivity & \textsc{Critical} &
  If Result B depends on Result A which depends on CY-A$_3$,
  then B is \emph{also} conditional on CY-A$_3$.
  Conditionality propagates through all downstream results.
  \textbf{Counter}: use \verb|\ClaimStatusConditional| with
  the full dependency chain.  Before writing
  \verb|\ClaimStatusProvedHere|, trace the proof chain back
  to its axioms. \\
\midrule
AP-CY12 & Shadow class from tower & \textsc{High} &
  G/L/C/M must be computed from the full shadow obstruction
  tower, not from generator counting or non-formality
  ($m_3 \neq 0$) alone.  Example: local $\bP^2$ is class M
  (infinite depth), \emph{not} class L.
  \textbf{Counter}: always compute the full tower before
  assigning a shadow class. \\
\midrule
AP-CY13 & Stale Part references & \textsc{Low} &
  After any Part restructuring, grep all three volumes for
  stale \texttt{Part\textasciitilde{}[IVXL]} references.
  7+ stale refs survived a single restructuring event.
  \textbf{Counter}: use \verb|\ref{part:...}| exclusively,
  never hardcode Part numbers. \\
\bottomrule
\end{longtable}
}%

%% =====================================================================
\section{Deep empirical anti-patterns: AP-CY14 through AP-CY20}
\label{sec:ap-cy14-20-full}

Identified through 100-commit deep archaeology.

{%
\small
\begin{longtable}{@{}p{1.6cm}p{2.4cm}lp{7.4cm}@{}}
\toprule
\textbf{ID} & \textbf{Name} & \textbf{Severity}
  & \textbf{Description and counter} \\
\midrule
\endhead
AP-CY14 & Unconstructed in thm & \textsc{Critical} &
  Any statement whose proof chain passes through $G(X)$
  at $d=3$, $A_{K3 \times E}$, or any unconstructed object
  \textbf{must} use \verb|\begin{conjecture}|, never
  \verb|\begin{theorem}| or \verb|\begin{proposition}|.
  The LLM pattern-matches on logical structure
  (``if $X$ then $Y$'') without checking whether $X$ exists.
  11+ instances fixed across 4 commits.
  \textbf{Counter}: default to \verb|\begin{conjecture}|
  in Vol~III. \\
\midrule
AP-CY15 & README inflation & \textsc{Medium} &
  README must not claim ``verified'' or ``proved'' for
  structural analogies or pattern matches.
  The README accumulates stronger claims than the \texttt{.tex}
  supports.
  \textbf{Counter}: after README edits, verify every
  ``proved''/``verified'' against corresponding
  \verb|\ClaimStatus| tags. \\
\midrule
AP-CY16 & Matrix size conflation & \textsc{Medium} &
  $\Sp_4$ quotient by $\pm I_4$ ($4 \times 4$), \emph{not}
  $\pm I_5$.
  $O(\Lambda^{3,2})$ quotient by $\pm I_5$ ($5 \times 5$).
  When two groups of different rank appear in the same formula,
  subscripts get harmonized to the most frequent.
  \textbf{Counter}: verify matrix dimensions match group rank. \\
\midrule
AP-CY17 & MF CY dimension & \textsc{High} &
  For $W \colon \bA^n \to \bA^1$, $\MF(W)$ is CY$_{n-2}$,
  \emph{not} CY$_{n-1}$ (Dyckerhoff).
  ADE in 2 variables: CY$_0$ (semisimple).
  Need 4 variables for CY$_2$, 5 for CY$_3$.
  \textbf{Counter}: verify $n - 2$ against desired CY dimension
  (see the table in Appendix~\ref{app:notation-conventions},
  Section~\ref{sec:nc-cy}). \\
\midrule
AP-CY18 & Lattice theta series & \textsc{Medium} &
  Leech theta: minimum norm${}^2 = 4$, first correction at
  $q^2$ not $q^1$.  The match with $1/\eta^{24}$ extends
  through $q^1$.
  Never conflate $j(\tau)$ coefficients with $V_\Lambda$
  character coefficients.
  \textbf{Counter}: verify by direct computation. \\
\midrule
AP-CY19 & $\hat{A}$-genus halving & \textsc{High} &
  $\hat{A}(x) = \frac{x/2}{\sinh(x/2)}$; convergence radius
  $= 2\pi$ (first pole of $\sin(x/2)$ at $x = 2\pi$).
  Dropping the $/2$ in the argument gives spurious radius
  $\pi$.  Appeared in 3+ independent computations.
  \textbf{Counter}: always include the $/2$ in the argument. \\
\midrule
AP-CY20 & Normal bdl $\neq$ spectral & \textsc{High} &
  The $\Z \times \Z$ grading from $N_{C/Y}$ connects to
  $(q,t)$ through the $\Omega$-background, \emph{not}
  directly.  The intermediary mechanism (equivariant
  localization, Nekrasov partition function, refinement) must
  be stated.
  \textbf{Counter}: name the intermediary; never write
  ``$N_{C/Y}$ grading $=$ $(q,t)$ parameters.'' \\
\bottomrule
\end{longtable}
}%

%% =====================================================================
\section{6d hCS session anti-patterns: AP-CY21 through AP-CY26}
\label{sec:ap-cy21-26-full}

Identified during the April 2026 holomorphic Chern--Simons session.

{%
\small
\begin{longtable}{@{}p{1.6cm}p{2.4cm}lp{7.4cm}@{}}
\toprule
\textbf{ID} & \textbf{Name} & \textbf{Severity}
  & \textbf{Description and counter} \\
\midrule
\endhead
AP-CY21 & $E_3$ bar class M & \textsc{High} &
  $(1+t)^{3g}$ holds for classes L and C only.
  Class M: infinite-dimensional ($d_4$ survives).
  Chain-level dimension for all classes: $P(q)^{3g}$.
  \textbf{Counter}: state the shadow class \emph{before}
  claiming $E_3$ bar cohomology. \\
\midrule
AP-CY22 & Miki is algebra-specific & \textsc{Medium} &
  The $S_3$ permutation of $(q_1, q_2, q_3)$ comes from the
  Weyl group of the CY torus, not from the $E_3$ operad in
  general.  Counterexample: $k[x]/(x^2)$ is $E_3$ but has
  no Miki automorphism.
  \textbf{Counter}: state it requires
  $U_{q,t}(\widehat{\widehat{\fgl}}_1)$ specifically. \\
\midrule
AP-CY23 & $\Eone$ not $\Einf$ bialgebra & \textsc{Critical} &
  The coproduct $\Delta_z$ lives on the $\Eone$ (ordered)
  side of the Swiss-cheese operad.  The $\Einf$ averaging
  map kills the Hopf structure:
  $\mathrm{av}(r(z)) = \kappa_{\mathrm{ch}}$.
  Li's vertex bialgebra framework ($\Einf$) is the wrong
  categorical home.
  \textbf{Counter}: formulate all Hopf data at the $\Eone$
  level using $B^{\mathrm{ord}}$ with deconcatenation. \\
\midrule
AP-CY24 & Docstring confabulation & \textsc{Medium} &
  Agents produce correct code but fabricate ``ground truth''
  values in docstrings.  The function computes correctly; the
  docstring claims wrong values for $n \ge 4$.
  \textbf{Counter}: verify every numerical value in docstrings
  against actual function output. \\
\midrule
AP-CY25 & $R$-matrix from vacuum & \textsc{High} &
  $R(z) = (\id \otimes S) \circ \Delta_z(1_A)$ is
  \textbf{wrong}: applying the coproduct to the vacuum and
  then the antipode yields $1 \otimes 1$ by the counit axiom.
  The correct $R$-matrix is characterized via the
  half-braiding $\sigma_A(z)$.
  \textbf{Counter}: never extract $R$ from $\Delta(1)$;
  always construct via the half-braiding. \\
\midrule
AP-CY26 & Verdier $\neq$ $\sigma_2$ inv. & \textsc{High} &
  $k^! = -k$ comes from Shapovalov form transposition
  (Verdier duality transposes the inner product), \emph{not}
  from $\sigma_2(-h_i) = -\sigma_2$ (false: $\sigma_2$ is
  degree-2 homogeneous, hence even under $h_i \to -h_i$).
  \textbf{Counter}: derive $k^!$ from Shapovalov/Verdier,
  not from $\sigma_2$ inversion. \\
\bottomrule
\end{longtable}
}%

%% =====================================================================
\section{Swarm-mined anti-patterns: AP-CY27 through AP-CY33}
\label{sec:ap-cy27-33-full}

Mined from 180-agent verification swarm, April 2026.

{%
\small
\begin{longtable}{@{}p{1.6cm}p{2.4cm}lp{7.4cm}@{}}
\toprule
\textbf{ID} & \textbf{Name} & \textbf{Severity}
  & \textbf{Description and counter} \\
\midrule
\endhead
AP-CY27 & Sandbox non-persist. & \textsc{High} &
  Background agents report successful file writes but files
  do \emph{not} persist to the main working tree (sandbox
  isolation).  Three engines were ``written'' and verified
  passing inside the sandbox but did not exist on disk.
  \textbf{Counter}: verify file existence with \texttt{ls}
  after agent completion. \\
\midrule
AP-CY28 & Pole-unsafe test pts & \textsc{High} &
  When testing $g(z)$ with poles at $z = \pm h_i$, test
  points must avoid these values.  For $h = (1,-2,1)$:
  poles at $z = \pm 1, \pm 2$; default $z = 2$ with
  $h_1 = 2$ gives $\varphi(2) = 0$, hence
  $g_{01}(2) = 1/0$.
  \textbf{Counter}: use $h = (37, 41, -78)$ for
  large-parameter safety. \\
\midrule
AP-CY29 & Wrong-repo writes & \textsc{Medium} &
  Agents write files to the wrong volume's directory.
  Example: an $\fsl_2$ Serre engine was written to
  \texttt{chiral-bar-cobar/compute/} (Vol~I) instead of
  \texttt{calabi-yau-quantum-groups/compute/} (Vol~III).
  \textbf{Counter}: verify the full path includes the
  correct repo root after any agent file write. \\
\midrule
AP-CY30 & Factored $\neq$ solved & \textsc{Critical} &
  $S_{ijk} = R_{ij} R_{ik} R_{jk}$ constructed from a
  YBE-satisfying $R$-matrix does \emph{not} automatically
  satisfy the Zamolodchikov tetrahedron equation.
  Proved: $O(\kappa_{\mathrm{ch}}^2)$ obstruction
  (thm:zte-failure).
  \textbf{Counter}: never assume pairwise consistency
  implies higher-order consistency. \\
\midrule
AP-CY31 & Spectral $z \neq$ ws $z$ & \textsc{High} &
  Drinfeld coproduct $\Delta_z$: Yangian spectral parameter
  (shift $u \to u - z$).
  OPE $T(z)T(w) \sim c/2 \cdot (z-w)^{-4}$: worldsheet
  insertion coordinate.  These are \emph{different} objects.
  Setting $z = 0$ in $\Delta_z$ removes the spectral shift
  (no OPE singularity).
  \textbf{Counter}: before any $z = 0$ argument, state
  whether $z$ is spectral or worldsheet
  (see Appendix~\ref{app:notation-conventions},
  Section~\ref{sec:nc-z}). \\
\midrule
AP-CY32 & Reorg.\ $\neq$ bypass & \textsc{Medium} &
  The 6d factorization homology route appears to bypass
  CY-A$_3$, but each subproblem (local $E_3$ algebra for
  compact targets, handle decomposition of K3, VOA
  identification of output) secretly requires the same
  chain-level data that CY-A$_3$ demands.
  The route \emph{reorganizes} the conjecture into
  subproblems but solves none independently.
  \textbf{Counter}: verify that every subproblem in an
  alternative route is resolved independently. \\
\midrule
AP-CY33 & Chain $\neq$ rational & \textsc{High} &
  $E_3$ structure is genuine at the \emph{chain} level but
  collapses to $\Etwo$ under Kontsevich formality (rational
  coefficients).  Physical content (Miki automorphism,
  factorization homology, tetrahedron corrections) lives at
  the chain level.  Formality destroys it.
  \textbf{Counter}: always state whether a claim about $\En$
  structure is at the chain level or the rational/formal
  level. \\
\bottomrule
\end{longtable}
}%

%% =====================================================================
\section{Cross-programme anti-patterns: AP150--AP157}
\label{sec:ap-cross-full}

These anti-patterns apply across all three volumes.

{%
\small
\begin{longtable}{@{}p{1.6cm}p{2.4cm}lp{7.4cm}@{}}
\toprule
\textbf{ID} & \textbf{Name} & \textbf{Severity}
  & \textbf{Description and counter} \\
\midrule
\endhead
AP150 & Confabulated composites & \textsc{Critical} &
  Agents stitch real ingredients (a categorical equivalence
  from paper A, a representation-theoretic identity from
  paper B) into composite structures that do not exist in
  the literature.  Each ingredient is real; the composite
  arrow has never been constructed.
  \textbf{Counter}: verify each arrow in a composite diagram
  independently.  If any arrow is unverified, the composite
  is conjectural. \\
\midrule
AP151 & $\hbar$ convention clash & \textsc{High} &
  Two definitions of $\hbar$ ($\hbar = \log q$ vs
  $\hbar = (\log q)/(2\pi i)$, or deformation parameter vs
  Planck constant) can coexist in one chapter when material
  is drawn from different sources.  The discrepancy cascades
  silently.
  \textbf{Counter}: one file, one $\hbar$.  Grep for
  existing definitions before introducing any $\hbar$.
  If a second convention is needed, use a distinct symbol
  with an explicit bridge identity. \\
\midrule
AP152 & ``Ordered'' ambiguity & \textsc{Medium} &
  ``Ordered product'' can mean
  (a) labeled/indexed (combinatorial, $\Eone$ bar),
  (b) time-ordered/radially-ordered (analytic, OPE), or
  (c) normally-ordered (Wick).
  These produce \emph{different} algebraic structures.
  \textbf{Counter}: bare ``ordered'' is forbidden.
  Always qualify: ``labeled-ordered,'' ``time-ordered,''
  or ``normally-ordered.'' \\
\midrule
AP153 & $E_3$ scope inflation & \textsc{High} &
  $E_3$ structure on Hochschild cochains (Deligne conjecture)
  requires the input to be $\Einf$ (commutative).  For
  $\Eone$ input, Hochschild cochains carry only $\Etwo$
  (Dunn: $\Eone \otimes \Etwo = E_3$, but $\Eone$ input
  contributes $\Eone$ not $\Einf$).
  \textbf{Counter}: verify the input algebra is $\Einf$
  before claiming $E_3$ structure on its Hochschild
  cochains. \\
\midrule
AP154 & Two $E_3$ structures & \textsc{Medium} &
  Even when $E_3$ is correctly invoked, two distinct $E_3$
  structures exist:
  (a) algebraic $E_3$ from Deligne's conjecture on
  $\HH^*(A)$ for $\Einf$ $A$, and
  (b) topological $E_3$ from framed little 3-disks on
  configuration spaces of $\R^3$.
  They agree under formality; differ at the chain level.
  \textbf{Counter}: specify which $E_3$ and whether
  formality is assumed. \\
\midrule
AP155 & Novelty overclaim & \textsc{Medium} &
  When $\Phi$ recovers a known invariant (Euler
  characteristic, Mukai vector, DT count) through a new
  construction, the \emph{invariant} is not new, only the
  \emph{construction path}.
  \textbf{Counter}: state ``$\Phi$ recovers the known
  invariant $X$ (originally due to [ref]) via a new
  construction path through the chiral bar complex.'' \\
\midrule
AP156 & Weierstrass $P_1$ ambig. & \textsc{Medium} &
  $\theta_1'/\theta_1$ (logarithmic derivative of $\theta_1$,
  quasi-periodic) vs Weierstrass $\zeta(z; \Lambda)$ differ
  by $\mathrm{Im}(z)$-dependent terms that matter for modular
  transformation.
  \textbf{Counter}: specify which convention
  ($\theta_1'/\theta_1$ vs Weierstrass $\zeta_\tau$) and
  state the quasi-periodicity explicitly. \\
\midrule
AP157 & Degeneration-type dep. & \textsc{High} &
  Different degenerations (large complex structure, conifold,
  orbifold, MUM, tropical) produce different chiral algebras
  with different $\kappa_\bullet$-values.  A statement proved
  ``in the degeneration limit'' is meaningless without
  specifying \emph{which} degeneration.
  \textbf{Counter}: name the degeneration type and state
  whether the result is specific to that type or holds for
  all degenerations. \\
\bottomrule
\end{longtable}
}%

%% =====================================================================
\section{Formula-mechanical: FM24}
\label{sec:ap-fm24-full}

{%
\small
\begin{longtable}{@{}p{1.6cm}p{2.4cm}lp{7.4cm}@{}}
\toprule
\textbf{ID} & \textbf{Name} & \textbf{Severity}
  & \textbf{Description and counter} \\
\midrule
\endhead
FM24 & B-cycle $i^2$ sign & \textsc{Critical} &
  In genus $\ge 1$, the B-cycle integral involves factors of
  $i$.  The error $i^2 = 1$ (instead of $i^2 = -1$)
  propagates silently and produces $|q| = 1$ instead of
  $|q| < 1$ for the nome, destroying convergence of all
  $q$-expansions.
  \textbf{Counter}: after any computation involving B-cycle
  integrals, verify that $|q| < 1$ (convergence of
  $q$-expansion).  If $|q| = 1$, trace back to an $i^2$ sign
  error.  Also verify $\mathrm{Im}(\tau) > 0$ is preserved
  by all transformations. \\
\bottomrule
\end{longtable}
}%

%% =====================================================================
\section{290-agent session anti-patterns: AP-CY35 through AP-CY52}
\label{sec:ap-cy35-52-full}

Identified during the April 2026 290-agent comprehensive session.

{%
\small
\begin{longtable}{@{}p{1.6cm}p{2.4cm}lp{7.4cm}@{}}
\toprule
\textbf{ID} & \textbf{Name} & \textbf{Severity}
  & \textbf{Description and counter} \\
\midrule
\endhead
AP-CY35 & $B^{(j)}$ hierarchy & \textsc{Critical} &
  $B^{(0)} = $ Connes $B$ (mixed complex).
  $B^{(j \ge 1)} = $ Connes \emph{hierarchy} ($\Sd$-framing).
  The mixed complex axiom $[b, B^{(0)}] = 0$ does \textbf{not}
  extend to $[b, B^{(j)}] = 0$.
  Three ``proofs'' were wrong because of this confusion.
  \textbf{Counter}: always specify which $B^{(j)}$ and never
  assume the mixed complex identity for $j \ge 1$. \\
\midrule
AP-CY36 & $\kappa_{\mathrm{ch}}$ wrong formula & \textsc{Critical} &
  $\sum (-1)^i \dim \HH_i$ gives $\chi_{\mathrm{top}}$
  ($= 24$ for K3), \textbf{not} $\kappa_{\mathrm{ch}}$ ($= 2$).
  The correct formula is the Hodge-filtered supertrace
  $\sum (-1)^q h^{0,q}$.  Serre duality kills non-$F^0$
  contributions.
  \textbf{Counter}: never compute $\kappa_{\mathrm{ch}}$ as
  alternating sum of $\HH_i$ dimensions.  Use
  $\mathrm{str}_{F^0}(q^{L_0})$. \\
\midrule
AP-CY37 & $\kappa_{\mathrm{BKM}}$ decomp.\ & \textsc{High} &
  $\kappa_{\mathrm{BKM}} = \kappa_{\mathrm{ch}} +
  \kappa_{\mathrm{cat}}$ is a \emph{coincidence} for $N = 1$
  (K3$\times E$).  The correct universal formula is
  $\kappa_{\mathrm{BKM}} = c_N(0)/2$ (Borcherds weight theorem).
  Fails for 7/8 diagonal Siegel orbifolds.
  \textbf{Counter}: use $c_N(0)/2$, never the naive decomposition. \\
\midrule
AP-CY38 & Class M $E_3$ bar $\neq \infty$ & \textsc{High} &
  Class M $E_3$ bar cohomology is $6^g$ (proved via K\"{u}nneth),
  \textbf{not} infinite-dimensional.  The $d_4$ differential
  kills $\Lambda^0$ and $\Lambda^3$, leaving $[0,3,3,0]$ at $g = 1$.
  Chain level is $P(q)^{6g}$; cohomology is $6^g$.
  \textbf{Counter}: state ``$6^g$ (closed form via K\"{u}nneth)''
  for class M, not ``infinite.'' \\
\midrule
AP-CY39 & Incompatibility thm & \textsc{Critical} &
  For single-object cyclic $A_\infty$ CY$_3$: $\mu_3 \neq 0$
  forces $\mu_2 = 0$ on the augmentation ideal.  Cross-arity
  cancellation is \textbf{impossible} at the naive level.
  The TCFT $B^{(2)}$ differs from naive pairwise contraction.
  \textbf{Counter}: never assume $\mu_2$ and $\mu_3$ can
  coexist on the same graded piece at the chain level. \\
\midrule
AP-CY40 & ProvedHere no proof & \textsc{Critical} &
  A theorem carrying \verb|\ClaimStatusProvedHere| \textbf{must}
  have a \verb|\begin{proof}| block.  The adversarial agent
  found \texttt{thm:cy-to-chiral-d3} had ProvedHere but no proof.
  \textbf{Counter}: grep for \verb|ProvedHere| and verify
  a \verb|\begin{proof}| block follows within 50 lines. \\
\midrule
AP-CY41 & Partial update contradiction & \textsc{High} &
  When upgrading a conjecture to theorem, \textbf{all} instances
  must be updated.  The session found $\sim$30 locations still
  saying ``open'' after CY-A$_3$ was proved.
  \textbf{Counter}: after any status change, grep all three
  volumes for the old status string and update every match. \\
\midrule
AP-CY42 & $\phi_{0,1}$ normalization & \textsc{High} &
  $c(-1) = 1$ (standard Gritsenko--Nikulin) vs $c(-1) = 2$
  (K3 elliptic genus $= 2\phi_{0,1}$).  The factor of 2 is
  $\kappa_{\mathrm{ch}}(\mathrm{K3})$.  Propagated silently
  across 3 engines.
  \textbf{Counter}: state which normalization convention
  is in force and verify against the K3 elliptic genus. \\
\midrule
AP-CY43 & Shadow--Feynman tautology & \textsc{Medium} &
  At $L \ge 4$, the Feynman engine calls the shadow recursion,
  making the match tautological.  Independent verification
  requires computing $m_k$ directly (e.g., from $k$-point
  conformal blocks).
  \textbf{Counter}: for $L \ge 4$, verify via an independent
  computation path, not through the shadow recursion. \\
\midrule
AP-CY44 & CY-D false at odd $d$ & \textsc{Critical} &
  $\kappa_{\mathrm{ch}} \neq \chi(\cO_X)$ when $d$ is odd,
  because Serre duality forces $\chi(\cO_X) = 0$ for all
  odd-dimensional CY, while $\kappa_{\mathrm{ch}}$ can be nonzero.
  Root cause: additivity vs multiplicativity.
  \textbf{Counter}: never write $\kappa_{\mathrm{ch}} =
  \chi(\cO_X)$ outside the scope $d = 2$, $h^{1,0} = 0$.
  Use the dimension-stratified formula. \\
\midrule
AP-CY45 & $N=2$ trivial braiding & \textsc{High} &
  At $N = 2$: $q^2 = 1$.  The double braiding $c_{V,W}^2 = \id$,
  giving a \emph{symmetric} (not modular) tensor category.
  Non-abelian MTC requires $N \ge 3$ where $q^2 \neq 1$.
  \textbf{Counter}: verify $q^2 \neq 1$ before claiming
  modular (non-symmetric) structure. \\
\midrule
AP-CY46 & No native CY$_4$ Yangian & \textsc{High} &
  $\pi_4(BU) = \Z$ obstructs $E_4$ structure.  The correct
  object for CY$_4$ is a $p_1$-twisted double current algebra.
  The E$_n$ cascade maximum is $E_3$ for \textbf{all} $d \ge 3$.
  \textbf{Counter}: never write ``$E_4$ Yangian'' or
  ``CY$_4$ Yangian.''  Use ``$p_1$-twisted double current
  algebra.'' \\
\midrule
AP-CY47 & Structure fn degree & \textsc{High} &
  Structure function degree comes from \emph{Mukai rank}
  ($= 24$ for K3), \textbf{not} Lie algebra dimension.
  For $E_8 \times E_8$: degree $(24, 24)$ from 24 Mukai
  directions, not $(500, 500)$ from $\dim(\frake_8) \times 2$.
  \textbf{Counter}: verify structure function degree against
  Mukai lattice rank. \\
\midrule
AP-CY48 & 3d$\to$6d lift rate & \textsc{Medium} &
  Only 24\% of 3d results lift to 6d.  Algebraic structures
  lift 100\%; topological structures lift 0\%.
  6d is \textbf{not} a dimensional upgrade of 3d.
  \textbf{Counter}: state the lift rate and specify
  which structures lift and which do not. \\
\midrule
AP-CY49 & Tautological tests & \textsc{High} &
  $\sim$10\% of agent-produced tests are tautological (testing
  hardcoded values against themselves).  Must verify via
  independent computation paths.
  \textbf{Counter}: every test must have at least two
  independent verification sources (AP10 protocol). \\
\midrule
AP-CY50 & Duplicate agent launches & \textsc{Medium} &
  When relaunching failed agents, check the agent registry
  to avoid running the same task twice.  Duplicate launches
  waste compute and create merge conflicts.
  \textbf{Counter}: check the agent registry before
  any relaunch.  Use unique task IDs. \\
\midrule
AP-CY51 & Rate-limited partial writes & \textsc{Medium} &
  When an agent is rate-limited, it may write engines and tests
  but \emph{not} manuscript material.  Relaunching from scratch
  duplicates the compute work.
  \textbf{Counter}: check disk for persisted files before
  relaunching.  Resume from the persisted state. \\
\midrule
AP-CY52 & Mega-file anti-pattern & \textsc{Medium} &
  Files exceeding 3000 lines should be split.
  \texttt{toroidal\_elliptic.tex} was 7190 lines;
  \texttt{k3\_times\_e.tex} was 5986 lines.  Both needed
  splitting into logical sub-chapters.
  \textbf{Counter}: when a \texttt{.tex} file exceeds 3000
  lines, split it by section.  Target 1000--2000 lines per file. \\
\bottomrule
\end{longtable}
}%

%% =====================================================================
\section{Decision trees}
\label{sec:ap-decision-trees}

The following decision trees codify the most frequently triggered
anti-patterns into operational templates.

\subsection{HZ3-1: Theorem vs.\ conjecture (AP-CY6, AP-CY14)}

Before writing \verb|\begin{theorem}| in Vol~III:
\begin{enumerate}[nosep]
\item Does the proof chain pass through $A_X$ for $d = 3$, $G(X)$,
      $C(\frakg,q)$, or any object whose existence is part of the
      $d = 3$ programme?
      \textbf{Yes} $\to$ \verb|\begin{conjecture}| +
      \verb|\ClaimStatusConjectured|.  \textbf{Stop.}
\item Does it pass through $A_X$ for $d = 2$ (CY-A$_2$ proved)?
      \textbf{Yes} $\to$ \verb|\begin{theorem}| or
      \verb|\begin{proposition}|; cite CY-A explicitly.
\item Pure categorical / VOA / Yangian statement (no functor
      invocation)?
      \textbf{Yes} $\to$ \verb|\begin{theorem}| or
      \verb|\begin{proposition}|; classical proof.
\item \textbf{Uncertain} $\to$ default \verb|\begin{conjecture}|.
      Downgrade is cheaper than retraction.
\end{enumerate}

\subsection{HZ3-2: $\kappa$-subscript (AP113)}

Before writing any $\kappa$:
\begin{enumerate}[nosep]
\item Subscript present?  Required: $\{\mathrm{ch}, \mathrm{cat},
      \mathrm{BKM}, \mathrm{fiber}\}$.
\item Forbidden subscripts: $\{\mathrm{global}, \mathrm{BPS},
      \mathrm{eff}, \mathrm{total}, \mathrm{naive},
      \mathrm{MacMahon}\}$.
      If you wrote BPS, you mean BKM.
\item For meta-naming (``$\kappa$-spectrum,'' ``$\kappa$-value''):
      write $\kappa_\bullet$ (the bullet denotes the indexing
      variable).
\end{enumerate}

\noindent
Routing (AP-CY55: distinguish manifold vs algebraization invariants):

\smallskip\noindent
\emph{Manifold invariants} (topological, independent of algebraization):
\begin{itemize}[nosep]
\item Holomorphic Euler char $\chi(\cO_X)$ $\to$
      $\kappa_{\mathrm{cat}}$.
\item Lattice rank / fiber structure $\to$
      $\kappa_{\mathrm{fiber}}$.
\end{itemize}

\noindent
\emph{Algebraization invariants} (depend on constructed algebra):
\begin{itemize}[nosep]
\item Chiral algebra $A_\cC$ / $\Phi(\cC)$ $\to$
      $\kappa_{\mathrm{ch}}$.
\item Borcherds--Kac--Moody / Igusa weight $\to$
      $\kappa_{\mathrm{BKM}}$.
\end{itemize}

\subsection{HZ3-3: Conditional propagation (AP-CY11)}

Before writing \verb|\ClaimStatusProvedHere|:
\begin{enumerate}[nosep]
\item Does this result's proof chain reach back to CY-A$_3$ or
      any unconstructed object?
      \textbf{No} $\to$ \verb|\ClaimStatusProvedHere| is valid.
      \textbf{Yes} $\to$ \verb|\ClaimStatusConditional| + name
      the dependency chain in the body of the statement.
\end{enumerate}

\subsection{HZ3-7: MF CY dimension (AP-CY17)}

Before any $\MF(W)$ CY claim:
\begin{enumerate}[nosep]
\item Identify $W \colon \bA^n \to \bA^1$.  What is $n$?
\item $\MF(W)$ is CY$_{n-2}$.  Check $n - 2$ against the desired
      CY dimension.
\item Verify: $n = 2$ (CY$_0$), $n = 3$ (CY$_1$),
      $n = 4$ (CY$_2$), $n = 5$ (CY$_3$).
\end{enumerate}

\subsection{HZ3-9: Shadow class assignment (AP-CY12)}

Before assigning a G/L/C/M class:
\begin{enumerate}[nosep]
\item Compute the full shadow tower $\{S_k\}_{k \ge 2}$.
\item $S_k = 0$ for all $k \ge 2$ $\to$ class G.
\item Tower terminates at finite depth $\to$ class L.
\item Tower terminates due to charge conservation $\to$ class C.
\item $S_k \neq 0$ for all $k$ (full tower computation required)
      $\to$ class M.
\item Generator counting or $m_3 \neq 0$ alone is \textbf{never}
      sufficient.
\end{enumerate}

%% =====================================================================
\section{Geometric vs Algebraic Model Conflations (AP-CY62--AP-CY67)}
\label{sec:ap-geom-alg-model}
%% Added 2026-04-16: from 15-agent adversarial swarm on two derived centers.

\subsection{AP-CY62: Geometric vs algebraic chiral Hochschild model}

Two chain-level models compute chiral Hochschild cochains:
\begin{enumerate}[label=(\alph*),nosep]
\item \emph{Geometric model} $C^\bullet_{\mathrm{ch,geom}}(\cA)$:
      sections over $\overline{C}_{n+1}(X)$ with log forms, three-component
      differential ($d_{\mathrm{int}} + d_{\mathrm{fact}} + d_{\mathrm{config}}$).
\item \emph{Algebraic model} $C^\bullet_{\mathrm{ch,alg}}(\cA)$:
      $\prod_{n \ge 0} \End^{\mathrm{ch}}_A(n{+}1)[-n]$ with
      $\delta f = [m, f]$ (Gerstenhaber bracket differential).
\end{enumerate}
These are quasi-isomorphic for logarithmic chiral algebras (via the
logarithmic comparison theorem).  At genus $\ge 1$, the geometric
model carries curve-dependent data (Green's functions, period
corrections) that the algebraic model lacks.

\emph{Counter.}  When chain-level structure matters (filtrations,
$\En$-structure, genus $> 0$), specify ``geometric (FM)'' or
``algebraic (bar/operadic).''  Bare $C^\bullet_{\mathrm{ch}}(\cA,\cA)$
is forbidden in chain-level arguments.

\emph{Status of comparison theorem.}  The bar complex comparison
(geometric $\leftrightarrow$ algebraic) is a named theorem
(\texttt{thm:geometric-equals-operadic-bar}).  The chiral Hochschild
comparison is only a remark (\texttt{rem:comparison-geometric-hoch},
\texttt{chiral\_center\_theorem.tex:346}).  This is a structural gap:
the comparison is invoked at $100{+}$ locations without a named theorem
backing it.

\subsection{AP-CY63: BD chiral operad vs algebraic End\textsuperscript{ch}}

Beilinson--Drinfeld define chiral operations via D-module maps:
$\End^{\mathrm{ch}}_M(n) = \Hom_{D(X^n)}(j_* j^* M^{\boxtimes n},
\Delta_* M)$.  The algebraic chiral endomorphism operad has
$\End^{\mathrm{ch}}_A(n) = \Hom(A^{\otimes n},
A((\lambda_1))\cdots((\lambda_{n-1})))$ (formal Laurent series).
These are isomorphic after formal-disk restriction and coordinate
choice (four-step chain: choose point, choose coordinate, trivialise
D-module, identify spectral variables with relative positions).

\emph{Counter.}  Never write ``the chiral endomorphism operad''
without specifying BD (D-module, coordinate-free) or algebraic
(formal Laurent series, coordinate-dependent).
The bridge is an isomorphism of non-$\Sigma$ operads,
valid after local coordinate trivialisation.  It requires the
$\Aut(\cO)$-equivariance statement.  A standalone Bridge
Proposition assembling the four steps is currently absent from
the manuscript.

\subsection{AP-CY64: Three-way Hochschild confusion (ChirHoch/HH/GF)}

Three invariants share the name ``Hochschild'' or ``derived centre'':
\begin{enumerate}[label=(\roman*),nosep]
\item $\ChirHoch^*(\cA)$: chiral Hochschild cochains.  Concentrated
      in $\{0,1,2\}$ (Theorem~H).
\item $\HH^*(A_{\mathrm{mode}})$: classical Hochschild cochains of
      the mode algebra (associative).  Concentrated in $\{0\}$ for
      simple algebras (Weyl algebra: $\dim = 1$).
\item $H^*_{\mathrm{GF}}(\mathrm{Lie}(A))$: Gel'fand--Fuchs
      continuous Lie cohomology.  Unbounded polynomial ring
      $\bC[\Theta_1, \ldots, \Theta_r]$.
\end{enumerate}
The claim ``Theorem~H has no THH analogue'' is \textbf{false} in
general: $\HH^*$ of the Weyl algebra is also concentrated (in
degree~$0$, even more so).  The genuine ``fails to concentrate''
object is (iii)~Gel'fand--Fuchs, which is neither ChirHoch nor~THH.

\emph{Escape route.}  At critical level $k = -h^\vee$,
$\ChirHoch^*$ becomes infinite-dimensional (Feigin--Frenkel centre
$\Fun(\Op_{\fg^\vee}(D))$), while $\HH^*(A_{\mathrm{mode}})$
remains finite.  This is the \textbf{only} regime where the two
derived centres genuinely differ in dimension.

\emph{Counter.}  Never claim ``ChirHoch* is finite while THH* is
infinite'' without specifying the level.  At generic level both are
finite.  At critical level, state the Feigin--Frenkel mechanism
explicitly.

\subsection{AP-CY65: Spectral parameter provenance}

The spectral parameter $z$ in the $R$-matrix $R(z)$ has a
\emph{three-part} origin:
\begin{enumerate}[label=(\alph*),nosep]
\item \emph{Algebraic}: the algebra $A$ must have a translation
      automorphism $\tau_z$ (creating evaluation modules $V_u$).
\item \emph{Geometric}: in the HT setup, $\tau_z$ is holomorphic
      translation on the curve~$C$ (the closed colour of $\SCchtop$).
\item \emph{Representation-theoretic}: the $R$-matrix depends on
      $z = u - v$ (evaluation parameter difference).
\end{enumerate}
The claim ``topological Drinfeld centre has no spectral parameters''
is \textbf{false}: the Yangian $Y(\fg)$, as a purely associative
algebra, has evaluation modules $V_u$ and spectral $R(z{=}u{-}v)$
in its Drinfeld centre.  The correct claim is about the bar complex:
the chiral bar \emph{differential} is $z$-dependent (OPE residues);
the topological bar \emph{coproduct} is $z$-independent
(deconcatenation).

\emph{Counter.}  Never write ``spectral parameters from the chiral
derived centre.''  Write: ``spectral parameters from the evaluation
module family, created by the translation automorphism that the
chiral structure provides.''

\subsection{AP-CY66: BZFN ambient category is not a tunable parameter}

The BZFN theorem (Lurie HA~5.3.1.30) states: for an $\Eone$-algebra
$A$ in a symmetric monoidal stable $\infty$-category~$\cS$,
$\cZ(\mathrm{LMod}_A(\cS)) \simeq
\mathrm{LMod}_{\HH^\bullet(A,A)}(\cS)$.  Both sides use the
\emph{same}~$\cS$.

The two derived centres arise from two \emph{different algebras}:
\begin{itemize}[nosep]
\item $A$ (chiral algebra in $\mathrm{IndCoh}(\Ran(X))$) $\to$
      $C^\bullet_{\mathrm{ch}}(A,A)$;
\item $A_{\mathrm{mode}}$ (mode algebra in $\Vect$) $\to$
      $\HH^\bullet(A_{\mathrm{mode}}, A_{\mathrm{mode}})$.
\end{itemize}

\emph{Counter.}  Never say ``applying BZFN in two different ambient
categories to the same algebra.''  Say: ``two different algebras
($A$ in D-modules, $A_{\mathrm{mode}}$ in $\Vect$), each with its
own BZFN equivalence.''

\subsection{AP-CY67: ``Spectral parameters from $\FM_k(\bC)$'' is narration}

The Vol~I preface (line~1014) says ``chiral Hochschild cochains
computed via the chiral endomorphism operad $\End^{\mathrm{ch}}_\cA$
with spectral parameters from $\FM_k(\bC)$.''  This compresses a
three-layer relationship into a prepositional phrase:
\begin{enumerate}[label=(\roman*),nosep]
\item \emph{Global geometric model}: sections over $\FM_{n+1}(X)$
      with log forms.
\item \emph{Formal-disk restriction}: choosing a local coordinate,
      the configuration space $C_k(\hat{D}_p)$ is parameterised by
      relative positions $\lambda_i = z_i - z_n$.
\item \emph{Algebraic model}: $\End^{\mathrm{ch}}_A$ with formal
      spectral variables $\lambda_i$.
\end{enumerate}
The spectral parameters \emph{do not come from} $\FM_k(\bC)$.
They are formal algebraic variables; their relationship to FM is
mediated by the local-global identification theorem (a comparison,
not a definition).

\emph{Counter.}  Replace ``spectral parameters from $\FM_k(\bC)$''
with ``spectral parameters from the chiral endomorphism operad,
identified with relative position coordinates on the formal disk via
the local-global comparison.''

\subsection{AP-CY68: Fiber-vs-total-space $\kappa_{\mathrm{cat}}$ for
  K3-fibered CY$_3$}

$\kappa_{\mathrm{cat}}(X) = \chi(\cO_X)$ is the holomorphic Euler
characteristic of the \emph{total space}.  By K\"unneth,
$\chi(\cO_{K3 \times E}) = \chi(\cO_{K3}) \cdot \chi(\cO_E)
= 2 \cdot 0 = 0$, so
$\kappa_{\mathrm{cat}}(K3 \times E) = 0$, not~$2$.
The value~$2$ is $\chi(\cO_{K3})$, the \emph{fiber} invariant.

This appears repeatedly in $\kappa_\bullet$-spectrum tables because
the conjectural BKM decomposition
$\kappa_{\mathrm{BKM}} = \kappa_{\mathrm{ch}} + \chi(\cO_{\mathrm{fiber}})$
uses the fiber~$\chi$ on the right-hand side; the temptation is to
relabel the fiber~$\chi$ as ``$\kappa_{\mathrm{cat}}$'' because it
arithmetically closes the $N{=}1$ coincidence
$5 = 3 + 2$.  Per AP-CY55, this conflates a \emph{total-space
manifold invariant} with a \emph{fiber invariant}; the two are
different per AP-CY55's manifold-vs-algebraization-invariant
discipline.  The numerical coincidence $5 = 3 + 2$ is real, but it
identifies $\kappa_{\mathrm{BKM}} - \kappa_{\mathrm{ch}}$ with
$\chi(\cO_{\mathrm{fiber}})$, not with $\kappa_{\mathrm{cat}}$ of
the total space.

The correct universal formula is $\kappa_{\mathrm{BKM}} = c_N(0)/2$
(Borcherds weight theorem, \verb|prop:bkm-weight-universal|), which
holds for all K3-fibered Class~A CY$_3$'s independently of any
fiber-total decomposition.

\emph{Counter.}  Every $\kappa_\bullet$-spectrum table listing
$\kappa_{\mathrm{cat}}(X)$ for $X = K3 \times E$ (or any K3-fibered
CY$_3$) must use the \emph{total-space} value $\chi(\cO_X)$, which
vanishes for odd-$d$ Serre manifolds.  If the fiber value is
relevant for a BKM decomposition argument, label it explicitly as
$\chi(\cO_{\mathrm{fiber}})$ or $\kappa_{\mathrm{cat}}(\mathrm{fiber})$,
never as bare $\kappa_{\mathrm{cat}}$.

\emph{Instances found.}  Preface L107 (fixed chunk~5);
\verb|quantum_chiral_algebras.tex| L870/L1020/L2453,
\verb|cy_holographic_datum_master.tex| L973,
\verb|en_factorization.tex| multiple sites (pending cross-file
cleanup).

\subsection{AP-CY69: Hochschild homology vs cohomology for Connes vs
  Gerstenhaber}

The $\mathbb{S}^d$-framing (Connes operator hierarchy) lives on
Hochschild \emph{homology} $\HH_\bullet(\cC)$.  The Gerstenhaber
bracket of degree~$1 - d$ lives on Hochschild \emph{cohomology}
$\HH^\bullet(\cC, \cC)$.  These are two different complexes, not
two structures on the same complex.  Writing ``Hochschild complex
$\HH_\bullet(\cC)$ receives the $\Sd$-framing and Gerstenhaber
bracket'' conflates them.

Why this matters: the Gerstenhaber bracket's shifted degree
$1 - d$ is the mechanism determining native $\En$ level in the
CY-to-chiral functor ($d = 2 \Rightarrow \Etwo$,
$d \ge 3 \Rightarrow \Eone$).  This degree count requires the
bracket to sit on cohomology (where the degree shift is
unambiguous); placing it on homology loses the shift and breaks
the $\En$-hierarchy argument.

\emph{Counter.}  Always specify: ``Hochschild homology
$\HH_\bullet(\cC)$ carries the $\Sd$-framing via the Connes
$B$-operator; Hochschild cohomology $\HH^\bullet(\cC, \cC)$ carries
the Gerstenhaber bracket of degree $1 - d$.''  Never attribute both
to a single unmarked ``Hochschild complex''.

\emph{Instances found.}  Preface L79 (fixed chunk~5).  Related to
AP-CY2 (CY trace target is $\HC^-_d$ not $\HH_d$) and to HZ3-5
(derived center vs Drinfeld center vs $\Etwo$).

\subsection{AP-CY70: Internal-development metadata in reader-facing
  prose}

Five species of manuscript-internal metadata were observed leaking
into preface prose during the Chriss--Ginzburg rectification pass:
\begin{enumerate}[label=(\alph*),nosep]
\item Session timestamps: ``(2026-04-17 inscription)'',
  ``The 2026-04-17 campaign closed\ldots'', ``pre-2026-04 phrasing
  superseded\ldots''.
\item Commit hashes: ``(commit \texttt{cade61c})'' as parenthetical
  explanation.
\item Manuscript-version self-reference: ``the central open problem
  of the first edition of this volume''.
\item Internal AP-tag citations in prose: ``Conjecture~CY-C,
  AP-CY60''.
\item Internal healing-status commentary: ``the original
  $\kappa_{\mathrm{ch}}$-trichotomy presentation is healed''.
\end{enumerate}
None of these have meaning for an external reader (referee,
\emph{Inventiones} reviewer).  They are artefacts of the
rectification audit trail, where intermediate states of the
manuscript are annotated for internal consistency.  Leaving them in
published prose produces a document that reads like a commit log
rather than a monograph.

\emph{Counter.}  Before shipping prose, grep for the species above:
\verb|2026-|, \verb|commit |, \verb|inscription|, \verb|campaign|,
\verb|AP-CY|, \verb|healed|, \verb|first edition|,
\verb|earlier phrasing|.  Every hit either (i)~strips the metadata
while preserving the mathematical claim verbatim, or (ii)~migrates
the content to a \verb|notes/| changelog file or a commit message.
The rule: a preface should not date itself.  Mathematical facts are
timeless; manuscript history is not.

\emph{Instances found and fixed.}  Preface L145, L419, L665, L356
(fixed chunks~4--5); preface L330 AP-tag (fixed chunk~7).

%% =====================================================================
\section{Cross-reference to Hot Zone}
\label{sec:ap-hot-zone-xref}

The ten most frequently triggered patterns (``Hot Zone'' entries
HZ3-1 through HZ3-10 in the development guide) correspond to the
following catalogue entries:

\begin{center}
\begin{tabular}{@{}lll@{}}
\toprule
\textbf{Hot Zone} & \textbf{AP IDs} & \textbf{Core issue} \\
\midrule
HZ3-1 & AP-CY6, AP-CY14 & Unconstructed object in
  \verb|\begin{theorem}| \\
HZ3-2 & AP113 & Bare $\kappa$ without subscript \\
HZ3-3 & AP-CY11 & Conditional propagation failure \\
HZ3-4 & AP-CY7 & CoHA $\neq$ $\Eone$-chiral algebra \\
HZ3-5 & AP-CY3, AP-CY4 & $\Etwo$ / Drinfeld / derived center
  conflation \\
HZ3-6 & AP-CY8 & Borcherds denominator $\neq$ bar Euler product \\
HZ3-7 & AP-CY17 & $\MF(W)$ CY dimension $n - 2$ not $n - 1$ \\
HZ3-8 & AP-CY10 & Flop $\neq$ Koszul dual \\
HZ3-9 & AP-CY12 & Shadow class from full tower, not shortcuts \\
HZ3-10 & AP-CY13 & Stale cross-volume Part references \\
\midrule
HZ3-11 & AP-CY35 & $B^{(j)}$ hierarchy confusion \\
HZ3-12 & AP-CY36 & $\kappa_{\mathrm{ch}}$ wrong formula \\
HZ3-13 & AP-CY40 & ProvedHere without proof block \\
HZ3-14 & AP-CY44 & CY-D false at odd $d$ \\
HZ3-15 & AP-CY49 & Agent tautological tests \\
\midrule
HZ3-16 & AP-CY62 & Geometric vs algebraic ChirHoch model \\
HZ3-17 & AP-CY63 & BD operad vs algebraic $\End^{\mathrm{ch}}$ \\
HZ3-18 & AP-CY64 & ChirHoch / HH / GF three-way confusion \\
HZ3-19 & AP-CY65 & Spectral parameter provenance \\
HZ3-20 & AP-CY66 & BZFN ambient category not tunable \\
HZ3-21 & AP-CY67 & ``from $\FM_k(\bC)$'' is narration \\
\bottomrule
\end{tabular}
\end{center}
